\section{Introduction} \label{sec:introduction}
  \gls{dram} is the dominant main memory technology used in contemporary computing 
  systems~\cite{paperOnDRAM}. \gls{dram} stores each bit of information in a memory cell that 
  consists of a capacitor and an access transistor. Due to charge leakage in the 
  capacitor, periodic refresh operations are required to maintain data integrity. 
  Otherwise, the stored data may be lost over time~\cite{paperTriggeringRHHardwareFault}.
  \\
  The need to accommodate ever-larger memory capacities in constrained physical space 
  has led to a continuous reduction in \gls{dram} cell dimensions and a corresponding increase 
  in memory density. Consequently, memory cells can function with lower operating voltages 
  and smaller charge levels, enabling higher storage capacity and improved energy 
  efficiency~\cite{paperAnotherFlip}. Unfortunately, this scaling led to a problem known as disturbance 
  error. A hardware-induced failure that arises from cell-to-cell interference in scaled 
  memory technologies~\cite{paperRHARetrospective}, leading in bit flips.
  \\
  Kim et al.~\cite{paperFlippingBits} exploited this error in 2014 as an attack, also known as Rowhammer. 
  Over the years, this attack was further developed, first on FPGA-based memory 
  controller, then x86 followed involuntarily, and later the attack also reached ARM 
  systems~\cite{paperRowhammeronaRPi}. Primarily on mobile phones and 
  chromebooks~\cite{paperDrammer,paperTRRespass,paperBlacksmith,paperDRAMA}. Ultimately, 
  it also found its way onto embedded systems and single-board computers. In recent years, 
  analyses have been carried out on some older Raspberry Pi models. However, due to 
  power and die area restrictions in LPDDR4~\cite{paperBlacksmith} as well as temperature controlled 
  refresh rate~\cite{thesisDRAMA, datasheetLPDDR4, docJEDECLPDDR4, docJEDECLPDDR4x}, this attack 
  remains relevant for contemporary Raspberry Pi models.
  \\
  While this bachelor thesis limits itself to the theory, analysis, and evaluation of 
  experiments, the practical work focuses on the implementation of a Rowhammer-based 
  proof-of-concept. This proof-of-concept is used to demonstrate the occurrence of bit flips on 
  modern Raspberry Pi platforms, specifically the Raspberry Pi 4B and Raspberry Pi 5.
  \\
  \\
  \textbf{Outline.} The remainder of this thesis is structured as follows. The background section 
  provides an overview of the technical foundations required for this work. The methodology 
  section describes the practical approach used to design, implement, and evaluate the 
  proof-of-concept. The lab setup section details the configuration of the test systems and 
  experimental environment. The attack primitives section introduces the fundamental techniques 
  required to construct the attack, followed by sections that describe the individual attack 
  primitives in terms of implementation, design choices, and practical constraints. 
  The conclusion synthesizes the key challenges and insights, and underscores the main findings of this work.
